% Generated draft: requires mathematical and source review.
\documentclass[11pt]{article}
\usepackage[margin=25mm]{geometry}
\usepackage{fontspec}
\setmainfont{Latin Modern Roman}
\usepackage{amsmath,amssymb,mathtools,unicode-math}
\setmathfont{Latin Modern Math}
\usepackage{enumitem,needspace,xurl,hyperref}
\setlength{\parindent}{0pt}
\setlength{\parskip}{5pt}
\setlength{\emergencystretch}{3em}
\newcommand{\N}{\mathbb{N}}
\newcommand{\Z}{\mathbb{Z}}
\newcommand{\Q}{\mathbb{Q}}
\newcommand{\R}{\mathbb{R}}
\newcommand{\C}{\mathbb{C}}
\newcommand{\F}{\mathbb{F}}
\newcommand{\A}{\mathbb{A}}
\newcommand{\PP}{\mathbb{P}}
\newcommand{\E}{\mathbb{E}}
\newcommand{\eps}{\varepsilon}
\newcommand{\dd}{\,\mathrm{d}}
\begin{document}
\section*{UnsolvedMath ID: 800229}
\textbf{Source problem code:} catalogue entry
\subsection{Definitions and mathematical statement}
The supplied source asks: ``Solve $\bar\partial$--Neumann in $L^2(d\mu)$, $d\mu$ a measure, even on the unit disc.'' Here the unit disc is
\[
\mathbb D=\{z\in\C:|z|<1\}.
\]
For a positive measure $\mu$ on a specified measurable space, $L^2(\mu)$ denotes the space of equivalence classes of measurable functions $f$ satisfying $\int |f|^2\,d\mu<\infty$. Positivity and the underlying measurable space are necessary specifications for this usual Hilbert-space interpretation; the source does not explicitly provide them.

A conditional operator formulation on $\mathbb D$ can be described as follows. Suppose suitable spaces $H_0=L^2(\mathbb D,\mu)$ and $H_1=L^2_{(0,1)}(\mathbb D,\mu)$ have been specified, together with a densely defined closed realization
\[
T:\operatorname{Dom}(T)\subset H_0\longrightarrow H_1
\]
of $\bar\partial$. Let $T^*$ be its Hilbert-space adjoint. In complex dimension one, the corresponding operator on $(0,1)$-forms is
\[
\Box_{\mu,1}=TT^*,\qquad
\operatorname{Dom}(\Box_{\mu,1})
=\{u\in\operatorname{Dom}(T^*):T^*u\in\operatorname{Dom}(T)\}.
\]
One possible solvability target is: for every
\[
f\in H_1\cap(\ker\Box_{\mu,1})^\perp,
\]
find $u\in\operatorname{Dom}(\Box_{\mu,1})\cap(\ker\Box_{\mu,1})^\perp$ such that $\Box_{\mu,1}u=f$. Whether a bounded solution operator or further estimates are required is unspecified.

\textbf{Unresolved scope.} This conditional formulation illustrates a possible meaning of the source request; the source does not select the form degree, operator realization, boundary conditions, admissible measures, or meaning of ``solve.'' For an arbitrary measure, even a distributional realization of $\bar\partial$ on $L^2(\mu)$ equivalence classes is not automatic. Consequently the supplied record does not determine a single precise quantified conjecture. The phrase ``even on the unit disc'' identifies a special case of interest but leaves the broader domain class unspecified.
\subsection{Short English statement}
The source asks for an $L^2(d\mu)$ theory of the $\bar\partial$--Neumann problem for a measure $\mu$, emphasizing the unit disc as an already interesting case. The admissible measures, operator domains, boundary conditions, form degree, and desired solvability or estimates are unspecified. The displayed operator formulation is conditional and requires human review.
\subsection{Sources}
\begin{itemize}
\item Local UnsolvedMath registry, permanent ID \texttt{800229}, exact source problem code \texttt{catalogue entry}; title: \emph{Several Complex Variables Problem List --- Problem 26}. The supplied record lists Buzzard, Fornaess, and Krantz under \texttt{proposed\_by}, with proposed year 1995.
\item Buzzard, Fornaess, Krantz, \emph{The Several Complex Variables Problem List} (1995), Problem 26. Supplied source URL: \url{https://arxiv.org/abs/math/9509202}. The provenance records access on 2026-10-16 and extraction by \texttt{arxiv\_latex}; these are supplied provenance claims, not independently verified here.
\item The local record labels this problem \texttt{open}. This supplied status is not an independently verified assessment of the literature. Supplied upstream revision: \texttt{d371ee917ac9647fbb06cd6fc529deec6f3dda82}.
\end{itemize}
\end{document}
